\documentclass[a4paper,10pt]{article}
\usepackage{fontawesome5}
\usepackage{graphicx, xcolor}
\definecolor{aqua}{RGB}{0,208,187}
\usepackage{geometry}
\usepackage{enumitem}
\usepackage{hyperref}
\usepackage{amsmath}
\geometry{left=1in, right=1in, top=1in, bottom=1in}
\setlength{\parindent}{0pt}

\begin{document}

\begin{center}
    {\LARGE \textbf{Ramez Abdalla, Ph.D.}}\\
    \vspace{2mm}
    \faMapMarker \, Machstraße 8-10, Vienna, Austria \\
    \faPhone \, +49 (0) 1520 363 9100 \\
    \faEnvelope \, \href{mailto:rmz.aziz@gmail.com}{ramez.maher.abdalla@gmail.com} \\
    \faLinkedin \, \href{    https://www.linkedin.com/in/ramez-abdalla-ph-d-8b1642b2/}
    {linkedin.com/in/ramez-abdalla-ph-d-8b1642b2} \\
    \faResearchgate \, \href{https://www.researchgate.net/profile/Ramez-Abdalla-5}{researchgate.net/profile/Ramez-Abdalla-5}
\end{center}

\vspace{2mm}

\section*{Summary}
Research scientist specializing in computational fluid dynamics, flow assurance, and artificial intelligence applications in energy systems. Expert in developing advanced numerical methods for multiphase flow simulation and pipeline network optimization. Primary research focuses on physics-informed machine learning for flow prediction, thermal-hydraulic analysis, and flow assurance challenges including wax deposition and hydrate formation. Published researcher with 5+ peer-reviewed journal articles in high-impact venues including the Journal of Fluid Mechanics. Pioneer in combining physics modeling approaches with modern machine learning techniques for enhanced accuracy in flow regime prediction and real-time monitoring of complex pipeline networks.

\section*{Research Profile}
\begin{description}[style=nextline, leftmargin=0.5cm]
    \item[Senior R\&D Engineer, Weatherford Scientific Computing and AI Team:] Principal investigator for physics-informed neural networks in reservoir simulation. Lead researcher in developing novel deep learning architectures for real-time production optimization. Published multiple peer-reviewed papers on AI applications in reservoir engineering.

    \item[Scientific Machine Learning Developer, Datagration:] Principal researcher in physics-constrained machine learning models for reservoir engineering. Developed patent-pending algorithms for real-time production optimization. Lead author of scientific publications on virtual flow metering systems.
    
    \item[Research \& Teaching Assistant, Clausthal University of Technology:] Applies advanced analytics and machine learning to optimize reservoir engineering and energy systems. Led the machine learning lab, contributing to peer-reviewed journals with research on innovative reservoir simulation algorithms.
     
     \item[Research \& Teaching Assistant, The American University in Cairo:] Conducts research on machine learning applications in reservoir and production engineering, guides practical implementations for enhanced operational efficiency.
     
    \item[Assistant Lecturer, Alexandria Institute of Engineering and Technology:] Instructs programming and database fundamentals, designs practical exercises enhancing engineering proficiency in software design and data management.   
\end{description}

\section*{Education}
\begin{itemize}[leftmargin=0.5cm]
    \item \textbf{PhD Degree, Subsurface Energy Systems}, Clausthal University, Niedersachsen, Germany \hfill 2021 - 2023 \\
    Dissertation: “Adaption of a Conceptual Data Analytics Framework for Subsurface Energy Systems: A Study of Predictive Maintenance, Virtual Flow Metering, and Production Optimization” \\
    %Focus: Developed a comprehensive data analytics framework tailored for subsurface energy systems. Conducted extensive research on predictive maintenance algorithms to anticipate equipment failures, implemented virtual flow metering techniques for real-time monitoring, and devised optimization strategies to enhance production efficiency. The work integrated advanced machine learning models and big data analytics to improve decision-making in energy systems.
    
    \item \textbf{Master's Degree, Petroleum Engineering}, Cairo University, Cairo, Egypt \hfill 2016 - 2018 \\
    Thesis: “Automatic Diagnosis of Sucker Rod Pumped Wells Using Machine Learning Algorithms” \\
    %Focus: Investigated the application of machine learning algorithms for the automatic diagnosis of sucker rod pumping systems. The research included the development and training of classification and regression models to detect anomalies and predict system failures. Emphasized the integration of data from various sensors and the creation of a diagnostic tool to assist engineers in maintaining optimal well performance.
    
    \item \textbf{Bachelor's Degree, Petroleum Engineering}, Faculty of Petroleum and Mining Engineering, Suez, Egypt \hfill 2009 – 2013 
    %Focus: Provided a strong foundation in petroleum engineering principles, including reservoir engineering, drilling engineering, production operations, and well logging. Engaged in hands-on training and projects that encompassed the full spectrum of petroleum exploration and production processes.
\end
{itemize}

\section*{Certificates}
\begin{itemize}[leftmargin=0.5cm]
    \item \textbf{Professional Certificate, Software Engineering}, Information Technology Institute (ITI), Alexandria, Egypt \hfill 2015 – 2016 
    %Focus: Completed a program in software engineering, covering key areas such as software development lifecycle, algorithms, data structures, and various programming languages. Specialized in both backend and frontend development, with practical projects utilizing technologies like Java, Ruby on Rails, Django, and React JS. The program included in-depth training in software design, testing, and deployment.
\end{itemize}


\section*{Employment}
\begin{description}[style=nextline, leftmargin=0.5cm]
 \item[Sep 2024 - Present] \textbf{Senior R\&D Engineer\\ Scientific Computing and AI Team}, Weatherford \\
    Responsibilities: Development of advanced computational models for fluid flow analysis and flow assurance in complex pipeline systems. Combine physics-based modeling with machine learning for enhanced prediction accuracy. 
    \item[Feb 2021 - Aug 2024] \textbf{Scientific Machine Learning Developer}, Datagration \\
    Responsibilities: Developing advanced data science tools for both backend and frontend applications in reservoir engineering within a digital data platform. This involves designing and implementing machine learning models to predict reservoir behavior, optimizing production strategies, and integrating these models with a scalable cloud-based infrastructure for real-time data analysis and visualization.
    
    \item[Feb 2021 - Dec 2023] \textbf{Research and Teaching Associate}, Technische Universität Clausthal\\
    Responsibilities: Supervising Bachelor and Master theses in the field of reservoir engineering, with a focus on the application of machine learning techniques. Leading the machine learning lab, conducting research on innovative algorithms for reservoir simulation, and publishing findings in peer-reviewed journals. Facilitating collaboration between academia and industry to advance the use of AI in reservoir management.
    
    \item[Jan 2019 - Feb 2021] \textbf{Research and Teaching Assistant}, The American University in Cairo \\
    Responsibilities: Conducting problem-solving sessions and laboratory demonstrations on reservoir fluid characterization, employing tools like Pipesim to simulate petroleum production systems. Developing and teaching curriculum that integrates machine learning applications in reservoir and production engineering, guiding students in practical implementations of AI techniques for optimizing production processes.
    
    \item[Jul 2018 – Jun 2019] \textbf{Petroleum Engineer and Data Scientist}, Myr:Conn Solutions GmbH \\
    Responsibilities: Monitoring and analyzing the performance of sucker rod pumping systems using data analytics. Developing predictive maintenance models to forecast equipment failures and optimize operational efficiency. Conducting statistical analyses and creating visualizations to support decision-making processes in the management of petroleum extraction operations.
    
    \item[Feb 2018 - Jun 2019] \textbf{Assistant Lecturer}, Alexandria Higher Institute for Engineering and Technology \\
    Responsibilities: Conducting laboratory sessions on statistical quality management, teaching programming (Programming I), and instructing on database fundamentals (SQL). Designing experiments and practical exercises that enhance students' understanding of quality control, software development principles, and data management techniques. Mentoring students in their final year projects, ensuring the application of theoretical knowledge to real-world engineering problems.
    
    %\item[Sep 2015 - Jun 2016] \textbf{Software Developer Trainee}, Information Technology Institute (ITI) \\
    %Responsibilities: Engaging in web development projects utilizing a range of software technologies. Backend development included Java Remote Method Invocation, Ruby on Rails, and Django on Python, focusing on creating robust and scalable server-side applications. Frontend development involved using React JS to build dynamic and responsive user interfaces. Participated in the complete software development lifecycle, from requirements analysis to deployment and maintenance.
\end{description}

\clearpage
\section*{Research Competencies}
\begin{itemize}[leftmargin=0.5cm]
    \item \textbf{Machine Learning Research:}
        \begin{itemize}
            \item Physics-Informed Neural Networks (PINNs) for reservoir simulation
            \item Deep Reinforcement Learning for production optimization
            \item Advanced time-series analysis and forecasting
            \item Uncertainty quantification in ML models
        \end{itemize}
    \item \textbf{Cloud AI/ML Infrastructure:}
        \begin{itemize}
            \item AWS SageMaker for model development and deployment
            \item AWS Bedrock for foundation models and generative AI
            \item MLOps pipeline design and implementation
            \item Distributed training and inference optimization
        \end{itemize}
    \item \textbf{Full-Stack Development:}
        \begin{itemize}
            \item Backend: FastAPI, Django, AWS Lambda, Step Functions
            \item Frontend: React, TypeScript, Chakra UI
            \item DevOps: Docker, Kubernetes, AWS ECS/EKS
            \item API Design: REST, WebSocket
        \end{itemize}
    \item \textbf{Scientific Computing \& Numerical Methods:}
        \begin{itemize}
            \item Finite Element Analysis for pipeline stress modeling
            \item Numerical methods for PDE solving (FVM, FDM)
            \item High-Performance Computing for flow simulations
        \end{itemize}
    \item \textbf{Flow Assurance \& Pipeline Systems:}
        \begin{itemize}
            \item Multiphase flow modeling and simulation
            \item Thermal-hydraulic analysis of pipeline networks
            \item Wax and hydrate formation prediction
            \item Flow regime characterization and optimization
        \end{itemize}
    \item \textbf{Domain Expertise:}
        \begin{itemize}
            \item Advanced fluid mechanics and rheology
            \item Transport phenomena in complex geometries
            \item Heat and mass transfer in pipeline systems
            \item Machine learning for flow pattern recognition
        \end{itemize}
\end{itemize}


\vspace{3mm}

\section*{Teaching Experience}
\begin{itemize}[leftmargin=0.5cm]
    \item \textbf{Reservoir Fluids Characterization} \\
    Providing in-depth knowledge of the properties and behaviors of reservoir fluids, crucial for reservoir engineering and management.

    \item \textbf{Fundamentals of Geo Production Systems} \\
    Exploring the principles and technologies involved in the production of geo-resources, focusing on system optimization using nodal analysis and a brief about artificial lift systems.

    \item \textbf{Offshore Production} \\
    Examining offshore production techniques, equipment, and challenges, with an emphasis on innovative solutions and safety protocols.

    \item \textbf{Directional Drilling} \\
    Teaching the technical aspects and applications of directional drilling, including trajectory planning and execution for optimized resource extraction.
    
    \item \textbf{Digital Drilling} \\
    Covering modern digital technologies in drilling operations, including real-time data analysis, automation, and predictive maintenance.

    \item \textbf{Fundamentals of Thermodynamics} \\
    Introducing core thermodynamics principles, including energy conservation, heat transfer, and system efficiency in engineering applications.

    \item \textbf{Database and Data Management Fundamentals} \\
    Covering database design, SQL, and data management techniques essential for efficient data handling in engineering applications.

    \item \textbf{Statistical Quality Management} \\
    Teaching statistical methods and quality control processes to ensure high standards in engineering projects and research.
\end{itemize}


\vspace{3mm}

\section*{Research Projects}
\begin{description}[style=nextline, leftmargin=0.5cm]
   

    \item[Cloud-Native Agentic AI APP for Reservoir and Production Engineering]
    \textbf{Research Outcomes:}
    \begin{itemize}
        \item Architected scalable ML platform using AWS SageMaker, Bedrock, and serverless technologies
        \item Implemented automated MLOps pipeline reducing model deployment time by 70\%
        \item Developed custom foundation models for reservoir analysis with 85\% accuracy
        \item Built full-stack application handling 10,000+ daily predictions
    \end{itemize}

    \item[Physics-Informed Neural Networks for Waterflooding Optimization] 
    \textbf{Research Outcomes:}
    \begin{itemize}
        \item Developed novel PINN architecture achieving 92\% accuracy in flow prediction
        \item Published findings in high-impact journal (Impact Factor: 4.2)
        \item Reduced computational time by 75\% compared to traditional methods
    \end{itemize}
    
    \item[Multi-Agent Reinforcement Learning for Enhanced Oil Recovery] Lead Researcher \\
    \textbf{Research Contributions:}
    \begin{itemize}
        \item Pioneered new MARL algorithm improving recovery efficiency by 15\%
        \item Successfully deployed in field trials with 30\% cost reduction
        \item Results presented at international conferences and published in peer-reviewed journals
    \end{itemize}
    
    \item[Autonomous Control of drilling Rigs (IOT)] Team Advisor Digital Drilling Lab at TU Clausthal with Clausthal University, Germany \\
    Reference: Prof. Dr.- Ing. Philip Jäger (Head of the Petroleum Production Systems Department at TU Clausthal), Dr. Ing Carlos Paz (Head of Digital Drilling Lab Technologies). \\
    Technologies: Django – Redis – MQTT – ZMQ – Websocket – JQuery.
    
    \item[Possible Applications of AI in Predictive maintenance in mature oilfields (Rühlermoor oilfield)] Exxon Mobil Cooperation with Clausthal University, Germany \\
    Reference: Prof. Dr.- Ing. Philip Jäger (Head of Petroleum Production Systems Department), Prof. Dr.- Ing Andreas Rauch (Head of Software Engineering Department), Dr. Ing Stefan Wittek, Dr. Ing Nelson Perozo.
    
    \item[Predictive Maintenance of the electrical submersible pumps using XGBoosting and PCA] Cairn Energy - Datagration \\
    Reference: Dipl- Ing Michael Stundner (Vice President of Technology).
    
    \item[Virtual Flow Metering on the electrical submersible pumps] Datagration LLC/ Clausthal University of Technology \\
    Reference: Prof. Philip Jaeger (Chair of Petroleum Production). \\
    Objective and Delivery: Model that can predict oil production for sucker rod pumps in X field.
    
\end{description}

\vspace{3mm}

% \section*{Technical Skills}
% \begin{itemize}[leftmargin=0.5cm]
%    \item \textbf{Programming Languages:} Python, R, C\#, JavaScript, Typescript, SQL
%    \item \textbf{Web Development:} Django, React, FastAPI
%    \item \textbf{Data Science and Machine Learning:} scikit-learn, TensorFlow, Keras, PyTorch, ONNX, Reinforcement Learning
%    \item \textbf{Reservoir Simulation:} Eclipse, Petrel, CMG
%    \item \textbf{Production Engineering Software:} Prosper, Pipesim, Q-Rod
%    \item \textbf{Database Management:} MySQL, PostgreSQL, MongoDB
%\end{itemize}

%\vspace{3mm}
\vspace{3mm}

%\clearpage
\section*{Technical Publications}
\subsection*{Journals}
\begin{enumerate}[leftmargin=0.5cm]
    \item Abdalla, Ramez, El Ela, Mahmoud Abu, and El-Banbi, Ahmed. Identification of Downhole in Sucker Rod Pumped Wells Using Deep Neural Networks and Genetic Algorithms, SPE Production \& Operations 35 (2020): 435–447. DOI: \href{https://doi.org/10.2118/200494-PA}{10.2118/200494-PA}.
    \item Abdalla, Ramez, Samara Hanin, Perozo Nelson, Carvajal Carlos, and Philip Jaeger. Machine Learning Approach for Predictive Maintenance of the Electrical Submersible Pumps (ESPs), ACS Omega 2022 7 (21), 17641-17651. DOI: \href{https://doi.org/10.1021/acsomega.1c05881}{10.1021/acsomega.1c05881}.
    \item Abdalla, Ramez, Hollstein, W., Carvajal, C.P. et al. Actor-critic reinforcement learning leads decision-making in energy systems optimization—steam injection optimization. Neural Computing \& Applications (2023). DOI: \href{https://doi.org/10.1007/s00521-023-08537-6}{10.1007/s00521-023-08537-6}.
    \item Abdalla, Ramez, Al-Hakimi Waleed, Perozo Nelson and Philip Jaeger. Real Time Liquid Rate and Water Cut Prediction from the Electrical Submersible Pump Sensors Data Using Machine-Learning Algorithms, ACS Omega 2023. DOI: \href{https://doi.org/10.1021/acsomega.2c07609}{10.1021/acsomega.2c07609}.
    \item Abdalla, Ramez, Nermine Agban, Christian Lüddeke, Dan Sui, Philip Jaeger. Multi-agent physics-informed reinforcement learning for waterflooding optimization, Franklin Open, Volume 10, 2025, 100229, ISSN 2773-1863. DOI: \href{https://doi.org/10.1016/j.fraope.2025.100229}{10.1016/j.fraope.2025.100229}.
\end{enumerate}


\subsection*{Book Chapters}
\begin{itemize}[leftmargin=0.5cm]
    \item Ramez Abdalla; Philip Jaeger; Chapter Data-Driven Virtual Flow Metering Systems, Book: Data Science and Machine Learning Applications in Subsurface Engineering; Edition First; Published 2024; Imprint CRC Press; Pages 17; eBook ISBN 9781003366980. DOI: \href{https://www.taylorfrancis.com/chapters/edit/10.1201/9781003366980-5/data-driven-virtual-flow-metering-systems-ramez-abdalla-philip-jaeger}{10.1201/9781003366980-5/data-driven-virtual-flow-metering-systems-ramez-abdalla-philip-jaeger}.
\end{itemize}

\subsection*{Conferences}
\begin{enumerate}[leftmargin=0.5cm]

    \item Abdalla, R., Gönczi, D., Toumi, O., Nikolaev, D., Manasipov, R., Stundner, M. 2024. Optimizing Waterflooding in Multi-Agent Systems through Decentralized Actor-Centralized Critic Reinforcement Learning. \textit{European Association of Geoscientists \& Engineers}, 2024(1), 1-5. DOI: \href{https://doi.org/10.3997/2214-4609.202439064}{10.3997/2214-4609.202439064}. Available at: \url{https://www.earthdoc.org/content/papers/10.3997/2214-4609.202439064}.

    \item Nikolaev, D., Hentges, L.D.O., Abdalla, R., Pechorskaya, E., Fuehrer, F., Justus, F. 2024. Automated Data Cleansing and Integrity for Enhanced Gas Production Analysis. \textit{European Association of Geoscientists \& Engineers}, 2024(1), 1-5. DOI: \href{https://doi.org/10.3997/2214-4609.202439061}{10.3997/2214-4609.202439061}. Available at: \url{https://www.earthdoc.org/content/papers/10.3997/2214-4609.202439061}.

    \item Manasipov, R., Nikolaev, D., Didenko, D., Abdalla, R., Stundner, M. 2023. Physics Informed Machine Learning for Production Forecast. In: \textit{Proceedings of the SPE Reservoir Characterisation and Simulation Conference and Exhibition}, Abu Dhabi, UAE, January 2023. DOI: \href{https://doi.org/10.2118/212666-MS}{10.2118/212666-MS}. Available at: \url{https://www.onepetro.org/conference-paper/SPE-212666-MS}.

    \item Abdalla, R., Nikolaev, D., Gönzi, D., Manasipov, R., Schweiger, A., Stundner, M. 2023. Deep Insight into Electrical Submersible Pump Maintenance: A Predictive Approach with Deep Learning. In: \textit{Proceedings of the SPE Offshore Europe Conference \& Exhibition}, Scotland, UK, September 2023. DOI: \href{https://doi.org/10.2118/215596-MS}{10.2118/215596-MS}. Available at: \url{https://www.onepetro.org/conference-paper/SPE-215596-MS}.

    \item Noshi, C.I., Rizk, M., Abdalla, R. 2019. An Intelligent Data-Driven Approach for Production Prediction. In: \textit{Proceedings of the Technology Conference Offshore}, Houston, Texas. OTC-29243-MS. DOI: \href{https://doi.org/10.4043/29243-MS}{10.4043/29243-MS}. Available at: \url{https://www.onepetro.org/conference-paper/OTC-29243-MS}.

    \item Reddicharla, N., Alnaqbi, K., Ali, M., Kumar, A., Gönczi, D., Abdalla, R. 2024. Machine Learning Driven Problem Detection for Production Optimization – Well Cease to Flow Prediction. In: \textit{Proceedings of the Abu Dhabi International Petroleum Exhibition and Conference (ADIPEC)}, Abu Dhabi, UAE. SPE-222548-MS. DOI: \href{https://doi.org/10.2118/222548-MS}{10.2118/222548-MS}. Available at: \url{https://www.onepetro.org/conference-paper/SPE-222548-MS}.

    \item Abdalla, R., Toumi, O., Gönczi, D., Sidaoui, A., Nikolaev, D., Schweiger, A., Schweiger, G. 2024. Physics-Informed Reinforcement Learning for Motor Frequency Optimization in Electrical Submersible Pumps: Enhancing AI-Led Decision-Making in Production Optimization. In: \textit{Proceedings of the SPE Artificial Lift Conference and Exhibition - Americas}, The Woodlands, Texas, USA. SPE-219542-MS. DOI: \href{https://doi.org/10.2118/219542-MS}{10.2118/219542-MS}. Available at: \url{https://www.onepetro.org/conference-paper/SPE-219542-MS}.
    \item Gönczi, D., Reddicharla, N., Rachapudi Venkata, S.R., Nikolaev, D., Abdalla, R., 2024. Machine Learning Driven Problem Detection for Production Optimization – Well Cease to Flow Prediction. In: \textit{Proceedings of the Fourth EAGE Digitalization Conference \& Exhibition}, Paris, France. DOI: \href{https://doi.org/10.3997/2214-4609.202439047}{10.3997/2214-4609.202439047}. Available at: \url{https://www.earthdoc.org/content/papers/10.3997/2214-4609.202439047}.


\end{enumerate}

%\section*{Languages}
%\begin{itemize}[leftmargin=0.5cm]
%    \item \textbf{Arabic:} Native
%    \item \textbf{English:} Fluent
%    \item \textbf{German:} Beginner 
%\end{itemize}

\vspace{3mm}

%\section*{References}
%\begin{itemize}[leftmargin=0.5cm]
%    \item \textbf{Prof. Dr. Philip Jaeger} \\
%    Chair of Petroleum Production, Technische Universität Clausthal \\
%    \textbf{Email:} \href{mailto:philip.jaeger@tu-clausthal.de}{philip.jaeger@tu-clausthal.de}

%    \item \textbf{Dr.-Ing. Carlos Paz} \\
%    Head of Digital Drilling Lab Technologies, Technische Universität Clausthal \\
%    \textbf{Email:} \href{mailto:carlos.paz@tu-clausthal.de}{carlos.paz@tu-clausthal.de}
    
%    \item \textbf{Dr.-Ing. Georg Pölzlbauer} \\
%     Head of Consulting Data Integration – finnova \\
%    \textbf{Email:} \href{mailto:georg.poelzlbauer@gmail.com}{georg.poelzlbauer@gmail.com}
%\end{itemize}

\end{document}
